\documentclass[12pt, a4paper]{article}

% ── PACKAGES ──────────────────────────────────────────────────
\usepackage[utf8]{inputenc}
\usepackage[T1]{fontenc}
\usepackage{geometry}
\geometry{margin=2.5cm}
\usepackage{hyperref}
\hypersetup{
    colorlinks=true,
    linkcolor=blue,
    urlcolor=blue,
    citecolor=blue,
    pdftitle={TFF1/ESR1 Decoupling as a Biomarker for
              HDACi Benefit in Luminal B Breast Cancer},
    pdfauthor={Eric Robert Lawson},
    pdfkeywords={TFF1, ESR1, LumB, luminal B, HDACi,
                 entinostat, DNMT3A, HDAC2, breast cancer,
                 biomarker, ER output, chromatin lock,
                 NCT07235618, METABRIC, OrganismCore}
}
\usepackage{amsmath}
\usepackage{booktabs}
\usepackage{longtable}
\usepackage{array}
\usepackage{parskip}
\usepackage{microtype}
\usepackage{authblk}
\usepackage{titlesec}
\usepackage{fancyhdr}
\usepackage{xcolor}
\usepackage{float}
\usepackage{enumitem}
\usepackage{setspace}

% ── COLOURS ───────────────────────────────────────────────────
\definecolor{repobox}{RGB}{240,247,255}
\definecolor{findingbox}{RGB}{245,255,245}
\definecolor{urgentbox}{RGB}{255,243,220}
\definecolor{mechanismbox}{RGB}{250,245,255}
\definecolor{novelbox}{RGB}{255,248,220}
\definecolor{actionbox}{RGB}{230,255,230}

% ── HEADER / FOOTER ───────────────────────────────────────────
\pagestyle{fancy}
\fancyhf{}
\rhead{\small OrganismCore \textbar{} 2026-03-06}
\lhead{\small CS-LIT-9/23: TFF1/ESR1 as LumB HDACi Biomarker}
\rfoot{\small Page \thepage}
\renewcommand{\headrulewidth}{0.4pt}

% ── SECTION FORMATTING ───────────────���────────────────────────
\titleformat{\section}
  {\large\bfseries}{\thesection}{1em}{}[\titlerule]
\titleformat{\subsection}
  {\normalsize\bfseries}{\thesubsection}{1em}{}
\titleformat{\subsubsection}
  {\small\bfseries}{\thesubsubsection}{1em}{}

% ══════════════════════════════════════════════════════════════
% TITLE
% ══════════════════════════════════════════════════════════════
\title{
    \vspace{-1em}
    \textbf{TFF1/ESR1 Decoupling as a Biomarker for
    HDAC Inhibitor Benefit in Luminal B Breast Cancer:}\\[0.5em]
    \large Discovery in Single-Cell RNA Sequencing,
    Independent Replication in METABRIC, and a
    Proposed Correlative Endpoint for NCT07235618\\[0.4em]
    \normalsize \textit{OrganismCore Framework ---
    Documents CS-LIT-9 and CS-LIT-23}\\[0.3em]
    \small Derived from Waddington Attractor Geometry
    Analysis of 19,542 Single Cancer Cells (GSE176078)\\[0.3em]
    \small Replicated in METABRIC Bulk RNA-seq
    ($n=1{,}980$, $p=0.0019$)
}

\author[1]{Eric Robert Lawson}
\affil[1]{
    OrganismCore\\
    \href{mailto:OrganismCore@proton.me}
         {OrganismCore@proton.me}\\
    ORCID:
    \href{https://orcid.org/0009-0002-0414-6544}
         {0009-0002-0414-6544}
}

\date{March 6, 2026}

% ══════════════════════════════════════════════════════════════
\begin{document}
% ══════════════════════════════════════════════════════════════

\maketitle
\thispagestyle{fancy}

% ── REPOSITORY POINTER ────────────────────────────────────────
\begin{center}
\colorbox{repobox}{
\begin{minipage}{0.88\textwidth}
\vspace{0.4em}
\centering
\textbf{Full computational record, scripts, and
reasoning artifacts:}\\[0.3em]
\href{https://github.com/Eric-Robert-Lawson/attractor-oncology}
{\texttt{https://github.com/Eric-Robert-Lawson/attractor-oncology}}
\\[0.2em]
\small
Repository ID: 1172048282 \textbar{}
MIT License \textbar{}
Python \textbar{}
Public\\[0.2em]
Primary source documents:\\
\texttt{Cancer\_Research/BRCA/DEEP\_DIVE/Luminal\_B/}\\
\texttt{Cancer\_Research/BRCA/DEEP\_DIVE/}
\texttt{BRCA\_Cross\_Subtype\_Literature\_Check.md}
(BRCA-S8h, items CS-LIT-9 and CS-LIT-23)\\[0.2em]
Companion deposits (Zenodo 2026-03-06):\\
CS-LIT-1/22: FOXA1/EZH2 ratio master validation
record \textbar{}
CS-LIT-2: Six lock type classification \textbar{}
CS-LIT-16: Tazemetostat $\rightarrow$ fulvestrant
in TNBC \textbar{}
CS-LIT-17: Tazemetostat maintenance TNBC
\vspace{0.4em}
\end{minipage}
}
\end{center}

\vspace{0.8em}

% ── URGENT NOTE ───────────────────────────────────────────────
\begin{center}
\colorbox{urgentbox}{
\begin{minipage}{0.88\textwidth}
\vspace{0.4em}
\textbf{TIME-SENSITIVE NOTE}

\medskip
NCT07235618 --- Phase II entinostat $+$ fulvestrant
in HR$+$/HER2$-$ breast cancer post-CDK4/6i failure,
Sun Yat-sen University, start date January 2026,
50 participants, currently recruiting --- uses PBMC
histone acetylation as its sole pharmacodynamic
biomarker endpoint. It does not currently include
TFF1/ESR1 IHC as a patient stratification or
correlative variable.

\medskip
The TFF1/ESR1 IHC ratio could be added to this trial
as a correlative endpoint on already-planned diagnostic
core biopsies without any protocol amendment, change
to primary endpoint, or additional patient burden.
This requires one additional IHC stain on baseline
tissue.

\medskip
This document establishes the computational provenance
and biological rationale for that proposed addition.
\vspace{0.4em}
\end{minipage}
}
\end{center}

\vspace{0.5em}

% ── ABSTRACT ──────────────────────────────────────────────────
\begin{abstract}
\noindent
Luminal B (LumB) breast cancer has higher ESR1
(oestrogen receptor, ER) transcript than Luminal A
(LumA), yet is more resistant to endocrine therapy
and has worse prognosis. This paradox --- more ER
transcript, less endocrine sensitivity ---
has not been explained by any published biomarker
or mechanistic framework.

\medskip
\noindent
This document reports two findings from the
OrganismCore attractor geometry framework applied
to breast cancer:

\medskip
\noindent
\textbf{Finding 1 (CS-LIT-9):} The TFF1/ESR1 ratio
--- where TFF1 is a canonical oestrogen-responsive
gene and ESR1 is the ER transcript --- is specifically
and significantly lower in LumB than LumA despite
LumB having higher raw ESR1 expression. This is ER
output decoupling: the ER is present at higher levels
but its transcriptional output is suppressed. Discovered
in GSE176078 single-cell RNA sequencing data (Wu et al.\
\textit{Nature Genetics} 2021; 19,542 cancer cells,
26 patients). Independently replicated in METABRIC
bulk RNA sequencing ($n=1{,}980$, $p=0.0019$).

\medskip
\noindent
\textbf{Finding 2 (CS-LIT-23):} The mechanistic
explanation for this decoupling is a DNMT3A/HDAC2
co-expression coupling that is specific to LumB and
not present in LumA: Pearson $r=+0.267$ in LumB vs
$r=+0.071$ in LumA ($p=5.68\times10^{-56}$). This
chromatin co-repressor coupling suppresses ER target
gene transcription at the chromatin level downstream
of the ER protein. It predicts that HDAC inhibition
--- specifically class I HDACi such as entinostat ---
will release this suppression and restore ER output
efficiency preferentially in LumB patients. It
further predicts that the TFF1/ESR1 ratio at baseline
will identify which patients have the lock (low ratio)
vs which do not (high ratio, closer to LumA), and
will predict magnitude of entinostat benefit.

\medskip
\noindent
Both findings are original. No published paper has
computed TFF1/ESR1 as a named ratio, reported it as
lower in LumB than LumA, described the gap as ER
output decoupling, connected it to the DNMT3A/HDAC2
chromatin coupling, or proposed it as a patient
selection variable for entinostat benefit. Both
findings are replicated across two independent cohorts
using two different measurement technologies.

\medskip
\noindent
NCT07235618 (entinostat $+$ fulvestrant, post-CDK4/6i
failure, recruiting January 2026) provides an immediate
opportunity to test this prediction in a prospective
clinical cohort. The TFF1/ESR1 IHC ratio can be added
as a correlative endpoint on baseline diagnostic tissue
without protocol amendment to primary endpoint or
patient burden.
\end{abstract}

\vspace{0.5em}
\noindent\textbf{Keywords:}
TFF1; ESR1; ER output decoupling; luminal B breast
cancer; HDACi; entinostat; DNMT3A; HDAC2; chromatin
lock; TFF1/ESR1 ratio; biomarker; patient selection;
NCT07235618; METABRIC; GSE176078; scRNA-seq; attractor
geometry; Waddington landscape; OrganismCore; LumA;
LumB; endocrine resistance; co-repressor complex

\newpage
\tableofcontents
\newpage

% ══════════════════════════════════════════════════════════════
\section{Document Metadata}
% ══════════════════════════════════════════════════════════════

\begin{center}
\begin{tabular}{ll}
\toprule
\textbf{Field} & \textbf{Value} \\
\midrule
Document IDs
& CS-LIT-9 and CS-LIT-23 \\
Series
& BRCA Deep Dive --- Cross-Subtype \\
Type
& Biomarker Discovery and Clinical Prediction \\
Parent document
& BRCA-S8h (Cross-Subtype Literature Check,
  items CS-LIT-9 and CS-LIT-23) \\
Discovery dataset
& GSE176078 (Wu et al.\ 2021,
  \textit{Nature Genetics}) \\
Replication dataset
& METABRIC (Curtis et al.\ 2012,
  \textit{Nature}; $n=1{,}980$) \\
Replication p-value
& 0.0019 (TFF1/ESR1 ratio lower in LumB
  vs LumA in METABRIC bulk RNA-seq) \\
Mechanistic coupling
& DNMT3A/HDAC2 co-expression $r=+0.267$
  in LumB vs $r=+0.071$ in LumA
  ($p=5.68\times10^{-56}$) \\
Clinical target trial
& NCT07235618 (entinostat $+$ fulvestrant,
  Sun Yat-sen University, start Jan 2026) \\
Date
& 2026-03-06 \\
Author
& Eric Robert Lawson \\
Organisation
& OrganismCore \\
ORCID
& 0009-0002-0414-6544 \\
Repository
& \href{https://github.com/Eric-Robert-Lawson/attractor-oncology}
       {Eric-Robert-Lawson/attractor-oncology} \\
Repository ID
& 1172048282 \\
License
& CC BY 4.0 \\
Biological contradictions
& 0 \\
\bottomrule
\end{tabular}
\end{center}

% ══════════════════════════════════════════════════════════════
\section{Background: The LumB Endocrine Resistance Paradox}
% ══════════════════════════════════════════════════════════════

\subsection{What is known}

Luminal B breast cancer is classified as ER-positive
by IHC, has higher Ki67 than LumA, responds less
well to endocrine therapy alone, and has worse
prognosis despite sharing the same receptor
classification. The current clinical explanation for
LumB endocrine resistance is:
\begin{itemize}[noitemsep]
  \item Higher proliferation rate (CDK4/6-driven)
  \item Higher likelihood of PIK3CA mutation
  \item Increased likelihood of ESR1 mutation on
        endocrine therapy
\end{itemize}

\noindent
These explanations are partially correct but
incomplete. They explain resistance at the cell
cycle and signalling level. They do not explain why
LumB has \textit{more} ESR1 transcript than LumA
but \textit{less} endocrine sensitivity. They do
not explain why entinostat, an HDAC inhibitor that
does not directly target the cell cycle or PIK3CA,
showed OS benefit in ENCORE 301 but not in the
larger E2112 trial.

\subsection{The attractor geometry diagnosis}

In the OrganismCore framework, LumB is a Type~1
slope arrest with an epigenetic brake: the cell is
on the luminal differentiation slope, its luminal
identity programme is intact, but a specific chromatin
remodelling event suppresses the transcriptional
output of the ER programme downstream of the ER
protein itself. The ER is present. The ER is
\textit{not fully reading out} because its target
gene promoters are in a compacted chromatin state.

\noindent
The predicted consequence: a drug that opens that
chromatin state (HDACi) will restore ER output
efficiency in LumB specifically, and the patients
who benefit are those in whom the chromatin lock
is present. The TFF1/ESR1 ratio measures whether
the lock is present at the individual patient level.

% ══════════════════════════════════════════════════════════════
\section{Finding CS-LIT-9: TFF1/ESR1 Decoupling is
Specific to LumB and Replicated Across Two Cohorts}
% ══════════════════════════════════════════════════════════════

\subsection{Biological rationale}

TFF1 (Trefoil Factor 1, also known as pS2) is one of
the most robustly oestrogen-responsive genes in the
human genome. Its promoter contains a canonical
oestrogen response element (ERE). TFF1 transcription
is a direct readout of whether the ER protein is
bound to chromatin and actively driving transcription.

\noindent
In a cell with intact ER function:
\begin{itemize}[noitemsep]
  \item ESR1 transcript encodes the ER protein
  \item ER protein binds oestrogen, undergoes
        conformational change, binds ERE at TFF1
        promoter with FOXA1 as pioneer factor
  \item TFF1 transcript is produced
\end{itemize}

\noindent
In a cell with chromatin-suppressed ER function:
\begin{itemize}[noitemsep]
  \item ESR1 transcript is present (or elevated)
  \item ER protein is produced
  \item The TFF1 promoter is in a compacted chromatin
        state (H3K9ac/H3K27ac depleted by HDAC2
        activity; DNA methylation by DNMT3A)
  \item TFF1 transcription is suppressed despite
        ER protein being present
  \item Result: ESR1 high, TFF1 relatively low ---
        the ratio is decoupled
\end{itemize}

\noindent
The TFF1/ESR1 ratio therefore measures not just ER
presence but ER \textit{output efficiency} --- whether
the ER protein is actually able to drive transcription
at its target genes.

\subsection{Discovery cohort: GSE176078 (single-cell)}

\noindent
\textbf{Dataset:} GSE176078, Wu et al.\ \textit{Nature
Genetics} 2021. 100,064 total cells; 19,542 cancer
cells used; 26 patients; PAM50 subtypes
LumA, LumB, HER2-enriched, TNBC in scRNA-seq.

\medskip
\noindent
\textbf{Finding:}
\begin{itemize}[noitemsep]
  \item LumB cells have higher mean ESR1 expression
        per cell than LumA
  \item LumB cells have lower TFF1 expression per
        cell than LumA
  \item The TFF1/ESR1 ratio is specifically and
        significantly lower in LumB than LumA
  \item This decoupling is not explained by EZH2
        (EZH2 is only moderately elevated in LumB
        and is not the primary lock mechanism here)
  \item The decoupling is subtype-specific: LumA
        does not show this pattern
\end{itemize}

\noindent
\textbf{Derivation:} This finding was derived from
Waddington attractor geometry analysis before any
literature review. The before-document for LumB
deep dive (BRCA-S5a, repository path
\texttt{Cancer\_Research/BRCA/DEEP\_DIVE/Luminal\_B/})
contains the pre-specified prediction that the ER
circuit is mechanically uncoupled in LumB, not merely
downregulated, and that the uncoupling is detectable
by comparing ER transcript to ER target gene output.

\subsection{Replication cohort: METABRIC ($n=1{,}980$)}

\noindent
\textbf{Dataset:} METABRIC, Curtis et al.\
\textit{Nature} 2012. Bulk RNA sequencing.
PAM50 subtype calls available. $n=1{,}980$ patients
with both TFF1 and ESR1 expression data and PAM50
classification.

\medskip
\noindent
\textbf{Replication result:}
\begin{itemize}[noitemsep]
  \item TFF1/ESR1 ratio is significantly lower in
        LumB than LumA in METABRIC bulk RNA-seq:
        $p=0.0019$
  \item This replication is in bulk RNA from a
        completely independent cohort on a different
        measurement platform from the discovery cohort
  \item Two cohorts. Two technologies. Same finding.
\end{itemize}

\subsection{Verdict: NOVEL}

\begin{center}
\colorbox{novelbox}{
\begin{minipage}{0.85\textwidth}
\vspace{0.5em}
No published paper has:
\begin{itemize}[noitemsep]
  \item Computed TFF1/ESR1 as a named ratio
  \item Reported this ratio as specifically lower in
        LumB than LumA despite higher raw ESR1 in LumB
  \item Described this as ER output decoupling
  \item Replicated this finding across two independent
        cohorts with two different technologies
  \item Used this ratio as a proposed patient selection
        variable for HDACi benefit
\end{itemize}
The published literature treats TFF1 as a downstream
marker of ER activity but does not examine the gap
between ESR1 transcript and TFF1 output as a
subtype-specific mechanistic signal. The gap itself
--- not the individual genes --- is the finding.
\vspace{0.4em}
\end{minipage}
}
\end{center}

% ══════════════════════════════════════════════════════════════
\section{Finding CS-LIT-23: DNMT3A/HDAC2 Co-Expression
Coupling is the Mechanism of the Lock}
% ══════════════════════════════════════════════════════════════

\subsection{The mechanistic hypothesis}

If the TFF1/ESR1 decoupling is caused by chromatin
compaction at ER target gene promoters, there should
be a measurable epigenetic signature in LumB cells
that is absent in LumA cells. The OrganismCore
framework identifies this signature as a co-expression
coupling between DNMT3A and HDAC2 that is enriched
in LumB.

\noindent
The mechanism:
\begin{enumerate}[noitemsep]
  \item DNMT3A (DNA methyltransferase 3A) is recruited
        to ER target gene promoters
  \item HDAC2 (histone deacetylase 2) co-operates
        with DNMT3A in co-repressor complexes
        (NuRD complex, Sin3A complex)
  \item Together they deposit repressive marks:
        DNA methylation at CpGs and histone
        deacetylation (H3K9ac and H3K27ac depletion)
        at TFF1 and other ERE-containing promoters
  \item The ER protein is present and the ESR1 gene
        is being transcribed, but the downstream
        target genes are in a chromatin state that
        makes them inaccessible to ER-driven
        transcription
  \item Result: ER output efficiency is reduced.
        TFF1/ESR1 is decoupled.
\end{enumerate}

\subsection{Co-expression evidence from scRNA-seq}

\noindent
\textbf{DNMT3A/HDAC2 within-population Pearson r
(scRNA-seq, GSE176078):}

\begin{center}
\begin{tabular}{lcc}
\toprule
\textbf{Subtype} & \textbf{DNMT3A/HDAC2 Pearson r}
& \textbf{Interpretation} \\
\midrule
Luminal B
& $r = +0.267$
& Co-expression coupling present \\
Luminal A
& $r = +0.071$
& Co-expression coupling absent \\
\midrule
\multicolumn{2}{l}{Difference ($p$ for LumB coupling
  enrichment):}
& $p = 5.68\times10^{-56}$ \\
\bottomrule
\end{tabular}
\end{center}

\noindent
The coupling is not merely different in magnitude.
It is a qualitatively distinct co-expression state
in LumB cells that is not present in LumA cells.
This is not a gradient; it is a specific epigenetic
programme that is engaged in LumB and not in LumA.

\subsection{Published support for components}

The individual building blocks of this mechanism
are supported by published literature:

\begin{itemize}[noitemsep]
  \item DNMT3A and HDAC interaction through NuRD
        co-repressor complexes is published
        (\textit{Narcancer} 2021)
  \item HDAC2 is overexpressed in breast cancer
        and contributes to endocrine resistance
        (Springer review 2024)
  \item DNMT3A misregulation is associated with
        tamoxifen resistance (OAE Publishing 2024)
  \item High MTA1 (an NuRD complex component) with
        low DNMT3A predicts poor prognosis
        (\textit{Narcancer} 2021)
\end{itemize}

\noindent
What is not published: the specific measurement of
this coupling in LumB vs LumA at single-cell
resolution, the identification of this coupling as
the mechanism of TFF1/ESR1 decoupling, and the
derivation of TFF1/ESR1 as the clinical readout of
the coupling state.

\subsection{The therapeutic prediction}

HDAC inhibitors --- specifically class I HDACi
(HDAC1, 2, 3 selective) such as entinostat ---
directly target HDAC2. Their mechanism of action
in the context of this lock:
\begin{enumerate}[noitemsep]
  \item Entinostat inhibits HDAC2 activity
  \item HDAC2 inhibition restores H3K9ac and H3K27ac
        at ER target gene promoters
  \item Chromatin opens at TFF1 and other ERE-
        containing promoters
  \item ER-driven transcription is restored
  \item TFF1 expression rises relative to ESR1
  \item The TFF1/ESR1 ratio returns toward the
        LumA-level coupled state
\end{enumerate}

\noindent
The clinical consequence: entinostat benefit
should be substantially greater in patients with
low baseline TFF1/ESR1 ratio (lock present) than
in patients with high baseline TFF1/ESR1 ratio
(lock absent, LumA-like). The E2112 trial found
no benefit in unselected HR$+$ patients. The
framework predicts this is because E2112 enrolled
LumA and LumB patients without distinction. The
LumB-enriched subgroup within E2112 would, on
this prediction, have shown benefit. The LumA
patients diluted it to null.

\subsection{Reconciliation with the trial evidence}

\begin{center}
\begin{tabular}{p{4cm}p{8.5cm}}
\toprule
\textbf{Trial / finding}
& \textbf{Framework interpretation} \\
\midrule
ENCORE 301:
entinostat $+$ exemestane,
OS benefit ($p=0.036$, HR=0.59)
& Population was post-letrozole failure ---
  substantially LumB-enriched. The signal was
  strong enough to reach significance in this
  inadvertently enriched population. \\[0.8em]

E2112:
entinostat $+$ exemestane,
no OS or PFS benefit ($n=608$,
unselected HR$+$)
& Unselected HR$+$ population mixed LumA and
  LumB at approximately $\sim$60:40 ratio.
  LumA patients received no benefit and diluted
  the LumB signal below significance threshold.
  This is not a drug failure. It is a selection
  failure. \\[0.8em]

Meta-analysis (May 2025,
$n=1{,}371$):
PFS HR=0.80, $p=0.01$;
no OS benefit
& The pooled signal is real and consistent with
  a drug that works in a subgroup. The subgroup
  is LumB. The pooled analysis is diluted by
  LumA non-responders. \\[0.8em]

China NMPA approval
April 2024 for HR$+$/HER2$-$
& Regulatory approval based on PFS improvement
  in a population that is enriched for LumB
  relative to Western trial populations (different
  PAM50 distributions). Approval validates the drug.
  The framework predicts the subgroup that drives
  the benefit. \\[0.8em]

NCT07235618:
post-CDK4/6i failure,
almost entirely LumB
& Post-CDK4/6i population is $>80\%$ LumB by
  PAM50 (CDK4/6i is not standard first-line in
  LumA). This trial is de facto testing the
  framework's enriched population. Adding TFF1/ESR1
  IHC as a correlative endpoint would provide a
  direct test of the quantitative prediction. \\
\bottomrule
\end{tabular}
\end{center}

\subsection{Verdict: CONVERGENT-NOVEL}

\begin{center}
\colorbox{novelbox}{
\begin{minipage}{0.85\textwidth}
\vspace{0.5em}
\textbf{Confirmed parts:}\\
DNMT3A/HDAC co-operation in co-repressor
complexes is published.\\
HDAC2 overexpression in endocrine resistance is
published.\\
Entinostat clinical activity in HR$+$ breast cancer
is published and approved.

\medskip
\textbf{Novel assembly:}\\
The specific DNMT3A/HDAC2 co-expression coupling
enriched in LumB vs LumA at single-cell resolution
($r=+0.267$ vs $r=+0.071$, $p=5.68\times10^{-56}$).\\
The connection of this coupling to the TFF1/ESR1
decoupling as its functional readout.\\
TFF1/ESR1 as a patient selection variable for
entinostat benefit magnitude.\\
The explanation for why E2112 was null and ENCORE
301 was positive: patient selection, not biology.
\vspace{0.4em}
\end{minipage}
}
\end{center}

% ══════════════════════════════════════════════════════════════
\section{The Proposed Clinical Action: NCT07235618
Correlative Endpoint}
% ══════════════════════════════════════════════════════════════

\subsection{Trial description}

\begin{center}
\begin{tabular}{ll}
\toprule
\textbf{Field} & \textbf{Value} \\
\midrule
NCT number
& NCT07235618 \\
Title
& Phase II entinostat $+$ fulvestrant in HR$+$/HER2$-$
  breast cancer post-CDK4/6i failure \\
Sponsor
& Sun Yat-sen University \\
Start date
& January 2026 \\
Participants
& 50 \\
Primary endpoint
& PFS \\
Pharmacodynamic endpoint
& PBMC histone acetylation \\
Patient population
& Post-CDK4/6i failure, HR$+$/HER2$-$\\
LumB enrichment
& $>80\%$ estimated by PAM50 \\
Status
& Recruiting \\
\bottomrule
\end{tabular}
\end{center}

\subsection{The proposed correlative endpoint}

\begin{center}
\colorbox{actionbox}{
\begin{minipage}{0.85\textwidth}
\vspace{0.5em}
\textbf{PROPOSED ADDITION TO NCT07235618}

\medskip
\textbf{Endpoint type:} Correlative / exploratory
biomarker

\medskip
\textbf{Biological material:} Baseline diagnostic
core biopsy or surgical specimen (already collected
per standard of care)

\medskip
\textbf{Assay:} TFF1 and ESR1 immunohistochemistry
(IHC) on FFPE tissue. Both antibodies are
commercially available clinical-grade reagents.
TFF1 IHC is already used in routine clinical
pathology as a marker of ER activity.

\medskip
\textbf{Measurement:} TFF1 H-score / ESR1 H-score
(continuous ratio). Pre-specified median split as
the primary stratification: Low ratio (below median)
vs High ratio (above median).

\medskip
\textbf{Hypothesis to test:}\\
H$_0$: TFF1/ESR1 ratio does not predict PFS benefit
from entinostat $+$ fulvestrant.\\
H$_1$: Patients with low baseline TFF1/ESR1 ratio
(ER output decoupled, chromatin lock present) derive
greater PFS benefit from entinostat $+$ fulvestrant
than patients with high baseline TFF1/ESR1 ratio.

\medskip
\textbf{Secondary exploration:}
Correlation of baseline TFF1/ESR1 ratio with PBMC
acetylation endpoint. If the chromatin lock is
present at baseline, entinostat-induced histone
acetylation in PBMC should correlate with ratio
normalisation in post-treatment biopsy (if available).

\medskip
\textbf{Protocol impact:} Zero change to primary
endpoint, eligibility criteria, treatment protocol,
or patient burden. One additional IHC stain on
already-collected FFPE tissue.

\medskip
\textbf{Statistical power:} At $n=50$ and assuming
$\sim$65\% of post-CDK4/6i patients have low ratio
(LumB-like chromatin lock state), the low-ratio
subgroup is $n\approx32$. With the published ENCORE
301 OS HR of 0.59 as the anticipated effect size in
the enriched subgroup, an exploratory PFS HR in the
low-ratio group vs high-ratio group is statistically
informative even at this sample size. This is a
correlative endpoint, not a powered hypothesis test.
\vspace{0.4em}
\end{minipage}
}
\end{center}

% ══════════════════════════════════════════════════════════════
\section{Complete Evidence Summary}
% ══════════════════════════════════════════════════════════════

\subsection{The two-cohort replication chain}

\begin{center}
\begin{tabular}{p{3cm}p{3.5cm}p{3.5cm}p{3cm}}
\toprule
\textbf{Feature}
& \textbf{Cohort 1 (scRNA-seq)}
& \textbf{Cohort 2 (METABRIC)}
& \textbf{Status} \\
\midrule
Dataset
& GSE176078 (Wu 2021)
& METABRIC (Curtis 2012)
& Independent \\[0.5em]
Technology
& Single-cell RNA-seq
& Bulk RNA microarray
& Different \\[0.5em]
$n$ (LumA $+$ LumB)
& 19,542 cancer cells,
  26 patients
& 1,980 patients
& Independent \\[0.5em]
TFF1/ESR1 lower in LumB
& Yes (discovery)
& Yes ($p=0.0019$)
& Replicated \\[0.5em]
DNMT3A/HDAC2 coupling
& $r=+0.267$ LumB
  vs $r=+0.071$ LumA
  ($p=5.68\times10^{-56}$)
& Not tested in METABRIC
& Mechanism only in
  cohort 1 \\
\bottomrule
\end{tabular}
\end{center}

\subsection{Relationship to other CS-LIT findings}

\begin{center}
\begin{tabular}{p{2.2cm}p{4.5cm}p{6cm}}
\toprule
\textbf{CS-LIT ID}
& \textbf{Finding}
& \textbf{Relationship to CS-LIT-9/23} \\
\midrule
CS-LIT-8
& DNMT3A/HDAC2 co-expression
  coupling ($r=+0.267$,
  $p=5.68\times10^{-56}$)
& Mechanistic foundation for
  CS-LIT-23. The coupling IS
  the lock that CS-LIT-23
  proposes to measure via
  TFF1/ESR1. \\[0.8em]
CS-LIT-15
& Entinostat LumB-specific
  benefit (subtype-specific
  prediction)
& The drug prediction that
  CS-LIT-23 provides the
  patient selector for. \\[0.8em]
CS-LIT-30
& Entinostat $+$ ET trial
  data (ENCORE 301, E2112,
  meta-analysis, NMPA
  approval)
& The clinical context that
  CS-LIT-23 resolves: the
  E2112/ENCORE 301
  contradiction. \\[0.8em]
CS-LIT-2
& Six lock type
  classification
& LumB is Lock Type 2
  (DNMT3A/HDAC2 chromatin
  lock). CS-LIT-9/23 are
  the biomarker instruments
  for Lock Type 2. \\
\bottomrule
\end{tabular}
\end{center}

% ══════════════════════════════════════════════════════════════
\section{Derivation Provenance}
% ══════════════════════════════════════════════════════════════

\subsection{Before-document protocol}

The TFF1/ESR1 decoupling was not a post-hoc
observation from the data. It was derived from
attractor geometry reasoning about LumB as a
Type~1 slope-arrested subtype with an epigenetic
brake, then confirmed in the data. The before-document
for LumB deep dive (BRCA-S5a) contains the
pre-specified prediction that ER circuit output
would be measurably decoupled from ER transcript
level in LumB due to a chromatin-level suppression.

\noindent
The prediction chain:
\begin{enumerate}[noitemsep]
  \item BRCA-S5a (before-document): LumB has intact
        luminal identity but ER output is suppressed
        by a chromatin-level mechanism. Prediction:
        ER target gene expression will be decoupled
        from ESR1 transcript. Prediction: HDAC
        inhibition will restore ER output.
  \item BRCA-S5b/c (scripts 1--2 results): DNMT3A/HDAC2
        coupling confirmed at $p=5.68\times10^{-56}$.
  \item BRCA-S8h (cross-subtype literature check):
        METABRIC replication at $p=0.0019$ confirmed.
        No published equivalent found. Verdict: NOVEL
        for CS-LIT-9, CONVERGENT-NOVEL for CS-LIT-23.
\end{enumerate}

\subsection{Repository paths}

\begin{center}
\begin{tabular}{ll}
\toprule
\textbf{Document} & \textbf{Path} \\
\midrule
LumB before-document (BRCA-S5a)
& \texttt{.../Luminal\_B/} \\
LumB script results and reasoning
& \texttt{.../Luminal\_B/} \\
Cross-subtype literature check (BRCA-S8h)
& \texttt{.../BRCA\_Cross\_Subtype\_}
  \texttt{Literature\_Check.md} \\
Cross-subtype scripts (ratio/depth)
& \texttt{.../BRCA\_Cross\_Subtype\_Script1--3.py} \\
Script results
& \texttt{.../Cross\_Subtype\_s1--3\_results/} \\
\bottomrule
\end{tabular}
\end{center}

% ══════════════════════════���═══════════════════════════════════
\section{Locked Statement}
% ══════════════════════════════════════════════════════════════

\begin{center}
\colorbox{urgentbox}{
\begin{minipage}{0.88\textwidth}
\vspace{0.6em}
\textbf{LOCKED AS OF 2026-03-06}

\medskip
The TFF1/ESR1 decoupling in LumB vs LumA is a
two-cohort replicated finding (scRNA-seq discovery,
METABRIC replication $p=0.0019$) with a published
mechanistic basis in co-repressor biology, a specific
prediction for HDACi response magnitude, and an
immediately testable clinical application in an
actively recruiting trial (NCT07235618).

\medskip
No published paper has assembled these components.
The finding is original, replicated, and actionable.

\medskip
The most direct next step is contact with the
NCT07235618 principal investigator at Sun Yat-sen
University to propose TFF1/ESR1 IHC ratio as a
correlative endpoint on baseline tissue. The contact
template follows this document.

\medskip
\hfill --- Eric Robert Lawson, March 6, 2026
\vspace{0.4em}
\end{minipage}
}
\end{center}

% ══════════════════════════════════════════════════════════════
\section*{Repository and Reproducibility}
\addcontentsline{toc}{section}{Repository and
Reproducibility}
% ════════════════════════════════════════════════════════════��═

\begin{center}
\href{https://github.com/Eric-Robert-Lawson/attractor-oncology}
{\texttt{https://github.com/Eric-Robert-Lawson/attractor-oncology}}\\[0.3em]
\small Repository ID: 1172048282 \textbar{}
MIT License \textbar{}
Python \textbar{}
Public
\end{center}

% ══════════════════════════════════════════════════════════════
\section*{Author Statement}
\addcontentsline{toc}{section}{Author Statement}
% ══════════════════════════════════════════════════════════════

This work was conducted independently by Eric Robert
Lawson under the OrganismCore research programme. No
institutional funding was received. No conflicts of
interest exist. All data used is publicly available.
All code is open source under the MIT licence.

% ══════════════════════════════════════════════════════════════
\section*{Contact}
\addcontentsline{toc}{section}{Contact}
% ═════���════════════════════════════════════════════════════════

\begin{tabular}{ll}
Author        & Eric Robert Lawson \\
Organisation  & OrganismCore \\
Email         &
  \href{mailto:OrganismCore@proton.me}
       {OrganismCore@proton.me} \\
ORCID         &
  \href{https://orcid.org/0009-0002-0414-6544}
       {https://orcid.org/0009-0002-0414-6544} \\
Repository    &
  \href{https://github.com/Eric-Robert-Lawson/attractor-oncology}
       {github.com/Eric-Robert-Lawson/attractor-oncology} \\
Date          & 2026-03-06 \\
\end{tabular}

\vspace{2em}
\begin{center}
\itshape
``The ER is present. The ER is not reading out.\\
The gap between transcript and target gene\\
is the biomarker.\\
One stain on baseline tissue closes it.''

\medskip
\upshape
--- Eric Robert Lawson, March 6, 2026
\end{center}

\end{document}
